\section{CONCLUSION}
This paper studied resilient restoration of power cyber-physical systems when the runtime knowledge used for recovery may itself be incomplete, stale, conflicting, replayed, or adversarially manipulated. We formulated this attack surface as \emph{recovery knowledge poisoning} and treated the recovery knowledge plane as an explicit security object rather than as trusted side information. KG-FedMARL separates a protected semantic basis from uncertain runtime assertions and limits two propagation paths: from adversarial evidence to recovery beliefs and action authority, and from contaminated recovery transitions to the deployed policy. The framework combines a provenance-aware temporal belief update with bounded source-group influence, a risk-aware recovery gate that separates uncertain semantic prerequisites from authoritative physical verification, and trust-gated online policy adaptation with quarantined replay and stable--candidate policy separation. Dynamic knowledge graphs, decentralized MARL, and federated training remain compatible implementation mechanisms, but none is assumed to provide the security guarantee by itself.

Under the stated trust boundary, the security analysis provides mechanism-level guarantees that connect the adversarial evidence budget and framework parameters to four quantities: accumulated recovery-belief distortion, false-authorization risk, post-attack belief-correction time, and cumulative deployed-policy deviation. Controlled experiments remain inside the corresponding sufficient envelopes. In the 20-seed IEEE 33-bus feeder study under medium mixed poisoning, KG-FedMARL reduces mean false authorization by 73.6\% relative to the latest-evidence baseline and by 90.0\% relative to the static-KG baseline, while its AURC remains within about 0.9\% of the oracle-information AC reference. The security ablations further separate the roles of provenance capping, forgetting, semantic authorization, and replay quarantine. On the 8500-node stress case, the conservative AC verifier rejects every additional restoration block from the configured initial state, producing a fail-safe but liveness-limited fallback regime. We report this limitation explicitly: the result supports bounded execution rather than successful large-feeder restoration and motivates feeder-aware candidate partitioning and verifier tuning.

These guarantees and experiments remain conditional on authenticated source identities, protected semantic and actuation components, correct provenance grouping and belief-update implementation, containment of the realized state in the verifier's conservative state set, and an uncompromised authoritative verifier. They do not establish global convergence of the underlying nonlinear MARL optimizer and do not cover arbitrary corruption of the trusted semantic basis, verifier code, protected actuation path, or all independent provenance groups. Within these boundaries, KG-FedMARL provides a concrete way to reason about and experimentally trace how far corrupted recovery knowledge can propagate through an online cyber-physical restoration loop.
